\documentclass[11pt,a4paper]{article}
\usepackage[top=1.7cm,bottom=1.5cm,left=1.5cm,right=1.5cm,headsep=8pt,footskip=16pt]{geometry}
\usepackage{enumitem}
\usepackage{amssymb}
\usepackage{array}
\usepackage[table]{xcolor}
\usepackage{fancyhdr}
\usepackage{lastpage}
\usepackage{multicol}
\usepackage{polyglossia}
\setdefaultlanguage[calendar=gregorian, numerals=maghrib]{arabic}
\setotherlanguage{english}
\IfFontExistsTF{Amiri}{%
  \newfontfamily\arabicfont[Script=Arabic, Scale=1.08]{Amiri}
}{%
  \newfontfamily\arabicfont[Script=Arabic, Scale=1.08, Path=../fonts/, Extension=.ttf,
    UprightFont=*-Regular, BoldFont=*-Bold, ItalicFont=*-Italic,
    BoldItalicFont=*-BoldItalic]{Amiri}
}
\IfFontExistsTF{AlBattar}{%
  \newfontfamily\khatt[Script=Arabic, Scale=1.15]{AlBattar}
}{%
  \newfontfamily\khatt[Script=Arabic, Scale=1.15, Path=../fonts/,
    BoldFont=ae_AlBattar.ttf]{ae_AlBattar.ttf}
}
\renewcommand{\headrulewidth}{0.8pt}
\renewcommand{\footrulewidth}{0.8pt}
\fancyhf{}
\rhead{\khatt المدرسة العليا للأساتذة ببوزريعة}
\lhead{\khatt ملحقة بومرداس}
\lfoot{\khatt الأستاذ: د. بن عبد الرزاق يوسف}
\cfoot{\thepage/\pageref{LastPage}}
\pagestyle{fancy}
\parindent 0pt
\setlength{\columnsep}{0.6cm}
\setlength{\columnseprule}{0.4pt}
\setlist[enumerate,1]{label=\textbf{\arabic*.}, leftmargin=*, itemsep=5pt, topsep=2pt}
\begin{document}

\linespread{1.08}\selectfont
\begin{center}
{\LARGE\khatt\textbf{الاختبار التاسع والعشرون في الذكاء الاصطناعي}}\\[4pt]
{\normalsize وضعيات إشكالية ــ أسئلة متعددة الاختيارات ــ الفصول من 1 إلى 6}
\end{center}
\vspace{-2pt}\rule{\textwidth}{0.6pt}\\[2pt]
\textbf{التعليمة:} يتضمن الاختبار خمس وضعيات إشكالية. اقرأ كل وضعية بعناية، ثم اختر التحليل أو التصرف الأنسب
من بين الاختيارات الخمسة؛ لكل وضعية إجابة صحيحة واحدة فقط.
اكتب \textbf{حرف} الإجابة الصحيحة في جدول الإجابات في آخر الاختبار.
لكل وضعية 4 نقاط (المجموع 20 نقطة).
\vspace{4pt}

\noindent
{\large\khatt\textbf{الوضعية 1}}\hfill{\small(الفصل 1)}\par\nopagebreak\vspace{3pt}
\setlength{\fboxsep}{7pt}\colorbox{gray!10}{\parbox{\dimexpr\linewidth-14pt}{طلب أستاذ من أداة توليدية أبياتًا على بحر الطويل لتمرين عروضي، فجاءت الأبيات فصيحة جميلة المعنى، لكن التقطيع أظهر كسرًا في عدد منها، وخلطًا بين تفعيلات الطويل والبسيط. وتريد زميلته استعمال الأبيات نفسها في اختبار. ويضيق الوقت قبل موعد الاختبار، وتريد الزميلة أن تتجنب إعادة كتابة الأبيات بنفسها. ما التقدير الأدق؟}}\par\nopagebreak
\begin{itemize}[leftmargin=1.0cm, labelsep=6pt, itemsep=3pt, topsep=4pt, parsep=0pt]
  \item[\textbf{أ)}] الأداة لا تستطيع إنتاج أي بيت موزون إطلاقًا، وكل الأبيات السليمة الوزن جاءت مصادفة، فلا يمكن الاستفادة منها حتى بعد التقطيع والتصحيح، ولا في أنشطة اكتشاف الكسر.
  \item[\textbf{ب)}] الأداة تنتج نصًا يشبه الشعر لغويًا، لكنها لا تضمن سلامة الوزن، لأن العروض يتطلب ضبطًا دقيقًا للمقاطع لا تتقنه دائمًا؛ فلا تُستعمل أبياتها في تمرين أو اختبار إلا بعد تقطيعها والتحقق منها.
  \item[\textbf{ج)}] الأداة تتقن العروض إتقانًا تامًا، والخلط بين الطويل والبسيط مقصود لتنويع التمرين، فتُستعمل الأبيات كما هي في الاختبار، لأن التنويع يرفع مستوى التحدي عند الطلبة.
  \item[\textbf{د)}] يكفي أن يطلب من الأداة «تأكد من الوزن» لتصبح الأبيات سليمة، ثم تُستعمل في الاختبار دون تقطيع، لأن الأداة تراجع تفعيلاتها بدقة عند الطلب الصريح في المحاولة الثانية، ويصبح التمرين جاهزًا في دقائق.
  \item[\textbf{هـ)}] الأبيات سليمة الوزن ما دامت فصيحة جميلة المعنى، لأن الفصاحة تستلزم سلامة الوزن في الشعر العربي، والكسر الظاهر خطأ في تقطيع الأستاذ، فتُستعمل في الاختبار مباشرة، فيقتصد الأستاذ وقت التقطيع.
\end{itemize}
\par\vspace{6pt}
\noindent
{\large\khatt\textbf{الوضعية 2}}\hfill{\small(الفصل 2)}\par\nopagebreak\vspace{3pt}
\setlength{\fboxsep}{7pt}\colorbox{gray!10}{\parbox{\dimexpr\linewidth-14pt}{لاحظت أستاذة أن أداة ذكاء اصطناعي تكمل بدقة أبياتًا مشهورة جدًا إذا ذُكر صدرها، لكنها حين طُلب منها إكمال بيت نادر لشاعر مغمور أكملته بعجز مختلف عن الأصل، دون أي تنبيه. ويريد التلاميذ استعمالها لحفظ نصوص المحفوظات. وستكون المحفوظات موضوع تقويم شفهي في نهاية الشهر. كيف تفسر ذلك؟}}\par\nopagebreak
\begin{itemize}[leftmargin=1.0cm, labelsep=6pt, itemsep=3pt, topsep=4pt, parsep=0pt]
  \item[\textbf{أ)}] الأداة أخطأت عمدًا في البيت النادر لتختبر انتباه المستخدم، لأن النماذج مصممة تعليميًا لإدخال أخطاء مقصودة في النصوص غير المشهورة، فيُنبَّه التلاميذ إلى ذلك فقط.
  \item[\textbf{ب)}] الأبيات المشهورة تكررت في البيانات، لكن النموذج ينبّه دائمًا عندما يولّد تكملة غير مؤكدة، فما دام لم ينبّه فالعجز صحيح، ويمكن للتلاميذ حفظه دون الرجوع إلى الكتاب، ويكسبون وقتًا في الحفظ.
  \item[\textbf{ج)}] الأداة تملك نسخة من كل الدواوين، والعجز المختلف رواية أخرى للبيت النادر لم تعرفها الأستاذة، فيُعتمد ما قدّمته الأداة في الحفظ لأنه مأخوذ من مصدر أوثق من الكتاب، فيحفظ التلاميذ الرواية الأوثق.
  \item[\textbf{د)}] الأداة تحفظ الأبيات المشهورة والنادرة بالدقة نفسها، والخطأ في البيت النادر عطل تقني عابر، فيكفي أن يعيد التلاميذ الطلب حتى يظهر العجز الصحيح ثم يحفظوه.
  \item[\textbf{هـ)}] الأبيات المشهورة تكررت كثيرًا في بيانات التدريب فرسخت أنماطها، أما النادرة فيولّد لها النموذج تكملة محتملة لا مسترجعة؛ فيُعتمد في الحفظ على النص المقرر لا على تكملة الأداة.
\end{itemize}
\par\vspace{6pt}
\noindent
{\large\khatt\textbf{الوضعية 3}}\hfill{\small(الفصل 3)}\par\nopagebreak\vspace{3pt}
\setlength{\fboxsep}{7pt}\colorbox{gray!10}{\parbox{\dimexpr\linewidth-14pt}{طلب أستاذ من الأداة «عرضًا تقديميًا عن المدرسة الرومانسية» لحصة مدتها عشرون دقيقة، فجاء نص من ثلاثين شريحة، في كل منها فقرة طويلة. ويريد عرضًا يناسب الوقت، تتضمن كل شريحة أفكارًا موجزة، مع ملاحظات يقولها شفهيًا. وقد اعتاد عرض الشرائح كما تولّدها الأداة دون تعديل. أي طلب أنسب؟}}\par\nopagebreak
\begin{itemize}[leftmargin=1.0cm, labelsep=6pt, itemsep=3pt, topsep=4pt, parsep=0pt]
  \item[\textbf{أ)}] «اختصر العرض قليلًا»، لأن الأداة تعرف المدة المناسبة لكل حصة، وتحدد عدد الشرائح وطول النقاط وملاحظات المتحدث تلقائيًا، دون حاجة إلى ذكر المدة أو الشكل.
  \item[\textbf{ب)}] «اكتب العرض في شريحة واحدة طويلة تتضمن كل شيء»، لأن الشريحة الواحدة تمنع تشتت التلاميذ، والأستاذ يقرأ النص كاملًا عليهم، فلا حاجة إلى نقاط موجزة أو ملاحظات شفهية، ويكون التركيز على كلام الأستاذ وحده.
  \item[\textbf{ج)}] «أعدّ ثماني شرائح بنقاط موجزة وملاحظات للمتحدث»، ثم يعرضها دون مراجعة المحتوى، لأن تحديد الشكل يضمن وحده صحة المعلومات والشواهد الشعرية الواردة في الشرائح.
  \item[\textbf{د)}] «أعدّ مخطط عرض من ثماني شرائح لحصة عشرين دقيقة: لكل شريحة عنوان وثلاث نقاط لا تتجاوز كل منها عشر كلمات، مع ملاحظات المتحدث أسفلها، وشاهد شعري واحد في شريحتين»، ثم يراجع المحتوى.
  \item[\textbf{هـ)}] «أعدّ ستين شريحة قصيرة جدًا»، لأن كثرة الشرائح تعني إيجازًا أكبر في كل منها، فتناسب العشرين دقيقة بالضرورة، دون حاجة إلى ملاحظات المتحدث أو تحديد الشواهد، ويستمتع التلاميذ بتتابع الصور السريع.
\end{itemize}
\par\vspace{6pt}
\noindent
{\large\khatt\textbf{الوضعية 4}}\hfill{\small(الفصل 4)}\par\nopagebreak\vspace{3pt}
\setlength{\fboxsep}{7pt}\colorbox{gray!10}{\parbox{\dimexpr\linewidth-14pt}{تريد أستاذة تسجيل حصصها بالصوت، ثم تفريغها وتلخيصها بأداة ذكاء اصطناعي عامة، لتوفير ملخصات للتلاميذ الغائبين. تتضمن الحصص مداخلات التلاميذ بأسمائهم، وأحيانًا أحاديث عن ظروف شخصية. ولم تُخبر التلاميذ ولا الأولياء ولا الإدارة بهذا المشروع. وترى أن الفكرة ستخدم التلاميذ كثيرًا، وتريد البدء في الأسبوع القادم. ما الإجراء الأسلم؟}}\par\nopagebreak
\begin{itemize}[leftmargin=1.0cm, labelsep=6pt, itemsep=3pt, topsep=4pt, parsep=0pt]
  \item[\textbf{أ)}] تحصل على إذن الإدارة وتُعلم التلاميذ والأولياء وفق تعليمات المؤسسة، وتحذف الأسماء والأحاديث الشخصية قبل رفع النص إلى الأداة، وتتحقق من شروط تخزينها، لأن أصوات القُصّر وأقوالهم معطيات شخصية.
  \item[\textbf{ب)}] تُعلم التلاميذ شفهيًا فقط، وترفع التسجيلات كاملة دون حذف، لأن علم التلاميذ يكفي قانونيًا وأخلاقيًا، ولا يلزم إشراك الأولياء أو الإدارة في نشاط داخل القسم، فيطمئن الجميع إلى سلامة النشاط.
  \item[\textbf{ج)}] تتخلى عن الملخصات وتمنع نفسها من أي تسجيل، لأن تسجيل الحصص ممنوع في كل الأحوال، حتى بإذن الإدارة وإعلام الأولياء وحذف المعطيات الشخصية قبل المعالجة.
  \item[\textbf{د)}] تسجّل وترفع التسجيلات كما هي، لأن الهدف تربوي نبيل، والأدوات العامة لا تحتفظ بالتسجيلات، فلا حاجة إلى إعلام أي طرف أو حذف الأسماء والأحاديث الشخصية.
  \item[\textbf{هـ)}] تحذف الأسماء فقط وتُبقي الأحاديث الشخصية، دون إعلام أحد، لأن حذف الأسماء يجعل التسجيل مجهول الهوية تمامًا، حتى لو بقيت تفاصيل يعرف بها الزملاء أصحابها، ويبقى النشاط مفيدًا ومحميًا.
\end{itemize}
\par\vspace{6pt}
\noindent
{\large\khatt\textbf{الوضعية 5}}\hfill{\small(الفصل 5)}\par\nopagebreak\vspace{3pt}
\setlength{\fboxsep}{7pt}\colorbox{gray!10}{\parbox{\dimexpr\linewidth-14pt}{وجد طالب للقصيدة نفسها روايتين مختلفتين في كلمتين، إحداهما في طبعة محققة والأخرى في موقع إلكتروني، فطلب من أداة ذكاء اصطناعي «الرواية الصحيحة». فاختارت رواية الموقع وعلّلت ذلك بأنها «الأشهر»، دون ذكر مصدر. ويريد أن يعتمد الرواية الصحيحة في مذكرته، ومشرفه متخصص في التحقيق. ما التصرف المنهجي الأنسب؟}}\par\nopagebreak
\begin{itemize}[leftmargin=1.0cm, labelsep=6pt, itemsep=3pt, topsep=4pt, parsep=0pt]
  \item[\textbf{أ)}] يعتمد الرواية التي يفضّلها أسلوبيًا، لأن اختلاف الروايات يعني أن النص مفتوح لاختيار القارئ، فلا معنى لتفضيل طبعة محققة على موقع إلكتروني في مثل هذه الحالات.
  \item[\textbf{ب)}] يسأل أداة ثانية، فإن اختارت الرواية نفسها اعتمدها، لأن اتفاق أداتين على رواية يرجّحها على ما في الطبعة المحققة، ولو لم تذكر الأداتان أي مصدر لاختيارهما، فيكون اختياره موضوعيًا.
  \item[\textbf{ج)}] لا يحسم الرواية بحكم الأداة، بل يعتمد الطبعة المحققة ويطّلع على ما ذكره المحقق من اختلاف الروايات ومصادرها، ويشير إلى الرواية الأخرى إن لزم، لأن الشهرة ليست معيار صحة الرواية.
  \item[\textbf{د)}] يحذف الكلمتين المختلف فيهما من القصيدة، ويكتب مكانهما نقاطًا، لأن أي اختيار بين روايتين لا يمكن تعليله علميًا، حتى بالرجوع إلى ما ذكره المحقق من مصادر الروايات.
  \item[\textbf{هـ)}] يعتمد رواية الموقع، لأن الأداة اطلعت على كل الروايات وحكمت بأنها الأشهر، والشهرة هي المعيار المعتمد في تحقيق النصوص التراثية، فلا حاجة إلى الرجوع إلى الطبعة المحققة، كما يفعل كثير من الباحثين.
\end{itemize}
\par\vspace{6pt}
\begin{center}
{\khatt\textbf{جدول الإجابات}}\\[3pt]
\renewcommand{\arraystretch}{1.6}
\begin{tabular}{|>{\centering\arraybackslash}m{1.9cm}|*{5}{>{\centering\arraybackslash}m{1.2cm}|}}
\hline
\textbf{رقم السؤال} & \textbf{1} & \textbf{2} & \textbf{3} & \textbf{4} & \textbf{5} \\ \hline
\textbf{الإجابة} &  &  &  &  &  \\[8pt] \hline
\end{tabular}

\vspace{8pt}
\renewcommand{\arraystretch}{2.2}
\begin{tabular}{|>{\centering\arraybackslash}m{7.5cm}|>{\centering\arraybackslash}m{3.2cm}|>{\centering\arraybackslash}m{3.2cm}|}
\hline
\textbf{الاسم واللقب} & \textbf{الفوج} & \textbf{العلامة} \\ \hline
 & & \hfill /20 \\ \hline
\end{tabular}
\end{center}
\end{document}
